
\chapter*{Preface}
\addcontentsline{toc}{chapter}{Preface} 
\markboth{Preface}{Preface}   

These notes grew out of my personal study of group representation theory. 
While standard textbooks (such as Serre's \emph{Linear Representations of Finite Groups}) 
provide an excellent foundation, I found that the transition from classical character theory 
to the modern module-theoretic and categorical perspectives was often scattered across 
different sources. 

\vspace{1em}
These notes are my personal attempt to organize this material into a single, 
coherent narrative. These notes are written for advanced undergraduates or beginning graduate students who have completed a first course in \textbf{abstract algebra and linear algebra}.

\vspace{1em}
The notes are organized into three logically distinct parts:

\noindent
\textbf{Chapters 1--9: Core Theory.} 
    These chapters cover the foundational material: the group algebra, linear representations, characters, Schur's lemma, orthogonality relations, the regular representation, and induced representations. This part assumes only basic linear algebra and group theory.

\noindent
\textbf{Chapters 10--11: Beyond the Semisimple Case.} 
    Chapter 10 deals with representations over characteristic $p$ fields, including the failure of semisimplicity, Brauer characters, and decomposition matrices. We only provide a first 
    glimpse and refer the reader to standard references (e.g., Alperin, 
    Curtis-Reiner) for the deeper modular representation theory. Chapter 11 explores representations over non-algebraically closed fields, covering Schur indices, Galois compatibility, and the theory of group characters over arbitrary fields.

\noindent
\textbf{Chapter 12: Applications and Outlook.} 
    This final chapter discusses applications to number theory (Artin $L$-functions and the Artin conjecture). 


In addition, Appendix A provides a self-contained introduction to module theory and category theory, which serves as the technical backbone for the main text. Appendix B contains the necessary algebraic number theory preliminaries.

\vspace{1em}
In the spirit of full transparency: 

\noindent
\textbf{The main text} (Chapters 1--12) were written and verified by me.

\noindent
\textbf{Appendices} were initially drafted with the assistance of a generative AI tool. I am currently in the process of rigorously verifying every definition, theorem, and proof in this appendix against standard references (such as Neukirch and Milne). Readers should treat this appendix with caution and are encouraged to report any inaccuracies.

\vspace{1em}
If you find any mathematical mistake, typo, or have a better proof, please open an Issue/Pull Request on my GitHub (\url{https://github.com/ZiyanQin0407/group-rep-notes.git})

\noindent
or Codeberg (\url{https://codeberg.org/Ziyan_Qin/group-rep-notes.git}) repositories, or reach out via email at qinziyan2028@outlook.com. I am deeply grateful for any corrections or suggestions.

\vspace{2em}
\begin{flushright}
    \textit{Ziyan Qin} \\
    \textit{\today}
\end{flushright}